% !TEX program = xelatex
% ============================================================
% 全国大学生数学建模竞赛（CUMCM）论文模板 — LaTeX 版
% 编译方式: xelatex paper.tex（连编两次）
% 说明:
%   - 电子提交时代论文第一页即为"题目+摘要+关键词"（保证书/编号页在报名系统内，
%     以当年官网《论文格式规范》为准）
%   - 字号: 正文小四(12pt)，章节标题按 黑体/宋体 排版（本模板已配置）
%   - 三线表用 booktabs；公式自动编号 (1)(2)(3)…；附录用 listings 贴代码
%   - 用 Word 的队伍请改用 paper-template.md 模板
% ============================================================
\documentclass[12pt,a4paper,fontset=windows]{ctexart}

% ---------- 页面与字号 ----------
\usepackage[top=2.54cm,bottom=2.54cm,left=3.17cm,right=3.17cm]{geometry}
\setlength{\parskip}{0.2em}

% ---------- 数学 ----------
\usepackage{amsmath,amssymb}

% ---------- 表格（三线表） ----------
\usepackage{booktabs}
\usepackage{multirow}
\renewcommand{\arraystretch}{1.3}

% ---------- 图片 ----------
\usepackage{graphicx}
\usepackage{float}
\usepackage{caption}
\captionsetup{font=small,labelfont=bf}

% ---------- 代码（附录） ----------
\usepackage{listings}
\usepackage{xcolor}
\definecolor{codegreen}{rgb}{0,0.5,0}
\definecolor{codegray}{rgb}{0.5,0.5,0.5}
\definecolor{codeblue}{rgb}{0.2,0.2,0.7}
\lstdefinestyle{matlabstyle}{
  language=Matlab, basicstyle=\ttfamily\small,
  keywordstyle=\color{codeblue}, commentstyle=\color{codegreen},
  stringstyle=\color{codegray}, numbers=left, numberstyle=\tiny\color{codegray},
  frame=single, breaklines=true, showstringspaces=false
}

% ---------- 超链接（打印稿无色） ----------
\usepackage[hidelinks]{hyperref}

% ---------- 章节标题样式 ----------
\ctexset{
  section = {
    format = {\heiti\zihao{4}},
    aftername = {},
    name = {},
    number = \chinese{section},
    beforeskip = 1.2em, afterskip = 1.0em
  },
  subsection = {
    format = {\heiti\zihao{-4}},
    name = {},
    beforeskip = 0.8em, afterskip = 0.6em
  }
}
% 手动控制: 正文用 \section*{一、问题重述}（编号已在标题文字里）

\linespread{1.3}

\begin{document}

% ============================================================
% 第一页: 题目 + 摘要 + 关键词
% ============================================================
\thispagestyle{plain}
\begin{center}
  {\heiti\zihao{3} 论文题目（三号黑体居中，可省略副题）}
\end{center}
\vspace{0.6em}

\noindent{\heiti\zihao{-4} 摘要}
\vspace{0.3em}

（摘要正文。建议 900-1200 字，五要素齐全：
①一句话点明研究背景与总任务；
②按"针对问题一/问题二/问题三/问题四"逐段写每问的模型、算法、关键数值结果；
③收尾一句模型评价与推广。摘要中必须出现每问的关键结果数值，
这是评委给分的第一来源。示例见 05-abstract-and-writing.md。）

\vspace{0.6em}
\noindent{\heiti 关键词：}（3-6 个方法名词，如: 时间序列分析 灰色预测 线性规划 蒙特卡洛 灵敏度分析）

\newpage

% ============================================================
% 正文（七章，章节标题格式: 一、问题重述）
% ============================================================

\section*{一、问题重述}

\subsection*{1.1 问题背景}
（用自己的话复述题目背景，1-2 段，可引用题目数据来源。）

\subsection*{1.2 问题重述}
（将题目要求改写为清晰的数学问题陈述，逐条列出需要解决的问题。）

\section*{二、问题分析}

\subsection*{2.1 总体思路}
（先讲清整体技术路线：数据预处理→相关性分析→建模→求解→检验→优化。
建议配一张总体思路流程图。）

\subsection*{2.2 问题一分析}
（分析难点与拟用方法。每问一小节。）

\section*{三、模型假设}

（4-6 条编号假设，每条必须合理且有依据，如: 假设1：未来一周内各品类进价稳定。）

\section*{四、符号说明}

\begin{table}[htbp]
\centering
\begin{tabular}{ccc}
\toprule
符号 & 含义 & 单位 \\
\midrule
$x_{it}$ & 第 $i$ 种单品第 $t$ 天的进货量 & kg \\
$p_{it}$ & 第 $i$ 种单品第 $t$ 天的定价 & 元/kg \\
$w_{it}$ & 第 $i$ 种单品第 $t$ 天的损耗率 & -- \\
\bottomrule
\end{tabular}
\end{table}

\section*{五、模型建立与求解}

\subsection*{5.1 问题一模型的建立与求解}

\subsubsection*{5.1.1 数据预处理}
（缺失值、异常值、重复值处理，日期连续性，分组聚合等。数据题必有此节。）

\subsubsection*{5.1.2 模型建立}
（模型推导、公式连续编号。示例公式:）
\begin{equation}
  \min z = \sum_{t=1}^{T}\sum_{i=1}^{I} c_{it} x_{it}
  \label{eq:objective}
\end{equation}

\subsubsection*{5.1.3 模型求解}
（求解算法、参数设置、运行环境。）

\subsubsection*{5.1.4 结果分析与检验}
（结果表 + 图 + 误差分析。每图后必须跟一句"由图X可看出……"结论。）

\subsection*{5.2 问题二模型的建立与求解}
（同 5.1 结构。）

\subsection*{5.3 问题三模型的建立与求解}
（同 5.1 结构。）

\subsection*{5.4 问题四模型的建立与求解}
（同 5.1 结构。）

\subsection*{5.5 灵敏度分析与模型稳定性}
（对关键参数做 $\pm 5\%$、$\pm 10\%$ 扰动，结果表 + 结论。
结论句式："结果表明，结果变化不超过 X\%，模型稳定性良好；其中某参数较敏感。"）

\section*{六、模型的评价、改进与推广}

（优点 2-3 条、缺点 2-3 条、改进方向、推广场景。示例:）

\noindent\textbf{优点：}（1）模型简洁，参数少，可解释性强；（2）双算法互验，结果可靠。

\noindent\textbf{缺点：}（1）未考虑……；（2）数据量的限制导致……

\noindent\textbf{改进与推广：}（1）……；（2）模型可推广至……领域。

\section*{参考文献}

% 书籍: 作者. 书名[M]. 出版地: 出版社, 出版年.
% 期刊: 作者. 论文名[J]. 杂志名, 卷(期): 起止页码, 出版年.
% 网页: 作者. 标题[EB/OL]. 网址, 访问时间.
[1] 姜启源, 谢金星, 叶俊. 数学模型（第五版）[M]. 北京: 高等教育出版社, 2018.

[2] 司守奎, 孙玺菁. 数学建模算法与应用（第二版）[M]. 北京: 国防工业出版社, 2015.

[3] 张三, 李四. 某问题研究[J]. 某某学报, 2020, 32(5): 12-20.

[4] 国家统计局. 中国统计年鉴 2023[EB/OL]. https://www.stats.gov.cn/, 2026-09-07.

% ============================================================
% 附录（代码清单 + 程序）
% ============================================================
\newpage
\appendix
\section*{附录}
（此处列代码清单: 文件名 + 功能一句话 + 运行环境，然后贴核心代码。）

\begin{lstlisting}[style=matlabstyle, caption={示例: main.m 主程序框架}, label={lst:main}]
% main.m - 问题一求解主程序
clear; clc;
data = readmatrix('data.xlsx');
fprintf('数据加载完成，共 %d 条记录\n', size(data, 1));
\end{lstlisting}

\end{document}
